% analyse_v6_v7.tex — Transition v6 → v7 : illinois_refine_adaptive
% Projet Riemann_Lab — hprzeta — 11 juin 2026
% Compilé avec : lualatex analyse_v6_v7.tex
\documentclass[11pt,a4paper]{article}

\usepackage{amsmath,amsfonts,amssymb}
\usepackage{geometry}
\usepackage{booktabs}
\usepackage{listings}
\usepackage{xcolor}
\usepackage{hyperref}
\usepackage{array}
\usepackage{microtype}

\geometry{margin=2.5cm}

\definecolor{codeblue}{rgb}{0.0,0.3,0.6}
\definecolor{codegray}{rgb}{0.4,0.4,0.4}
\definecolor{codegreen}{rgb}{0.0,0.5,0.2}
\definecolor{backcolor}{rgb}{0.97,0.97,0.97}

\lstset{
  backgroundcolor=\color{backcolor},
  basicstyle=\ttfamily\small,
  commentstyle=\color{codegray}\itshape,
  keywordstyle=\color{codeblue}\bfseries,
  stringstyle=\color{codegreen},
  breaklines=true,
  frame=single,
  rulecolor=\color{codegray!40},
  language=C,
  numbers=none,
  xleftmargin=4pt,
  xrightmargin=4pt,
}

\hypersetup{colorlinks=true, linkcolor=codeblue, urlcolor=codeblue}

\title{\textbf{Analyse v6 $\to$ v7 : \texttt{illinois\_refine\_adaptive}}\\[4pt]
  \large Optimisation 2-phases de l'affinage Illinois en C/libmpfr}
\author{hprzeta --- Riemann\_Lab}
\date{11 juin 2026}

\begin{document}
\maketitle

\begin{abstract}
\noindent
La version v6 du pipeline de calcul des zéros de $\zeta$ identifie l'affinage
Illinois (\texttt{illinois\_refine}, \textsc{libmpfr} 170 bits) comme bottleneck à
\textbf{78--83\,\% du temps CPU}.
La v7 introduit \texttt{illinois\_refine\_adaptive}, un affinage en deux phases :
les premières itérations utilisent une précision réduite de 64~bits (1 limbe \textsc{mpfr}),
avant de basculer à 116~bits pour la convergence finale.
Le gain mesuré est \textbf{$\times 3.7$ sur la durée totale} à $T = 100\,000$
(30,9 min vs $\approx$113 min pour v6).
\end{abstract}

%----------------------------------------------------------------------
\section{Rappel v6 --- état et bottleneck}
%----------------------------------------------------------------------

La v6 (\textit{commit \texttt{b676e88}}) a résolu le problème des zéros manquants
introduit par un STEP adaptatif insuffisant : avec \texttt{STEP = 0.010} fixe et la
détection \texttt{scan\_arb.c} (Z double en C pur), le run $T = 100\,000$ donne
\textbf{138\,069 zéros, 0 manquant, Turing--Backlund COMPLET}.

Le profil de phases de v6 révèle un goulot unique :

\begin{center}
\begin{tabular}{lrrr}
\toprule
Phase & Temps cumulé & ms/appel & \% mur$\times$W \\
\midrule
\texttt{illinois\_C} & $\approx$5\,500 s & $\approx$130 ms & \textbf{83\,\%} \\
turing & $\approx$60 s & --- & 1\,\% \\
detection & $\approx$30 s & --- & 0,5\,\% \\
\bottomrule
\end{tabular}
\end{center}

L'affinage de chaque zéro par \texttt{illinois\_refine} (170 bits, 15--20 itérations,
$N_\text{full}$ termes par appel) est linéaire en
\begin{equation}
  N_\text{full}(t) = \left\lfloor\sqrt{\frac{t}{2\pi}}\right\rfloor
  \quad\Longrightarrow\quad
  N_\text{full}(100\,000) \approx 126\ \text{termes.}
\end{equation}
Avec 170 bits (3 limbes GMP), chaque opération \texttt{mpfr\_cos}, \texttt{mpfr\_log}
est lente. C'est la cible de v7.

%----------------------------------------------------------------------
\section{Principe de l'optimisation 2-phases}
%----------------------------------------------------------------------

L'algorithme Illinois sur $[a, b]$ converge en $\approx 15$--20 itérations pour
$|b - a| = \texttt{STEP} = 0.01 \to \varepsilon = 10^{-12}$.
Les 8 premières itérations réduisent l'intervalle de $0.01$ à $\approx 10^{-4}$
(gain $\approx 10^{8}$ sur la largeur) : c'est là que réside le coût dominant.

\textbf{Idée v7 :} évaluer $Z(t)$ en précision réduite (64 bits) durant la phase de
réduction rapide, puis en précision complète (116 bits) pour la convergence finale.

\begin{center}
\renewcommand{\arraystretch}{1.3}
\begin{tabular}{lcccl}
\toprule
Phase & Itérations & Précision & Limbes mpfr & Objectif \\
\midrule
1 (rapide) & $i = 0 \dots 7$ & 64 bits & 1 & $|b-a| \to 10^{-4}$ \\
2 (précise) & $i \ge 8$ & 116 bits & 2 & $|b-a| \to 10^{-12}$ \\
\bottomrule
\end{tabular}
\end{center}

La transition de phase est opérée par \texttt{mpfr\_prec\_round} qui élargit la
représentation sans recalcul.

%----------------------------------------------------------------------
\section{Décision technique : $N_\text{full}$ et non $N_\text{fast}$}
%----------------------------------------------------------------------

L'idée initiale était d'utiliser $N_\text{fast} = \max(5,\, N_\text{full}/4)$ termes
en phase 1, réduisant le coût d'évaluation d'un facteur $\times 4$.
\textbf{Cette approche a été abandonnée} pour la raison suivante.

La formule de Riemann--Siegel
\begin{equation}
  Z(t) = 2\sum_{n=1}^{N} \frac{\cos\!\left(\theta(t) - t\ln n\right)}{\sqrt{n}}
         + R(t), \quad N = N_\text{full}(t)
\end{equation}
n'est correcte que pour $N = \lfloor\sqrt{t/2\pi}\rfloor$.
Tronquer à $N_\text{fast} < N_\text{full}$ ne produit pas une approximation
grossière du même $Z$ : les termes manquants peuvent \emph{changer le signe} de la somme.

\begin{center}
\small\textit{Exemple} : $t = 1002$, $N_\text{full} = 12$, $N_\text{fast} = 3$.\\
$Z_\text{RS}(c,\, N=3)$ peut avoir le signe opposé à $Z_\text{RS}(c,\, N=12)$.
\end{center}

Un signe incorrect dans l'algorithme Illinois viole l'invariant
$Z(a) \cdot Z(b) < 0$ et la convergence diverge vers un pseudo-zéro arbitraire.
\textbf{La solution retenue conserve $N_\text{full}$ termes dans les deux phases} :
le gain de vitesse vient uniquement de la précision réduite.

\subsection*{Pourquoi 64 bits $\gg$ 170 bits en vitesse ?}

\textsc{mpfr} stocke un nombre flottant en $\lceil p/64 \rceil$ limbes de 64 bits.
Pour $p = 64$ bits : 1 limbe. Pour $p = 170$ bits : 3 limbes.
Les fonctions transcendantes (\texttt{mpfr\_cos}, \texttt{mpfr\_log}) exploitent des
routines SIMD spécialisées pour 1 limbe, absentes à 3 limbes.
Le rapport de vitesse observé est \textbf{$\times 10$--$20$} par appel, bien supérieur
au simple rapport de limbes ($\times 3$).

%----------------------------------------------------------------------
\section{Implémentation}
%----------------------------------------------------------------------

Deux fonctions ont été ajoutées dans \texttt{illinois\_mpfr.c} (commit \texttt{8637098}) :

\begin{lstlisting}[language=C, caption=Signature de \texttt{Z\_rs\_mpfr\_ntermes}]
/* Z(t) RS avec N_termes impose (pas floor(sqrt(t/2pi))) */
double Z_rs_mpfr_ntermes(double t_d, int N_termes, mpfr_prec_t prec);
\end{lstlisting}

\begin{lstlisting}[language=C, caption=Coeur de \texttt{illinois\_refine\_adaptive}]
double illinois_refine_adaptive(
    double a, double b, double fa, double fb,
    double t, int iter_switch, int max_iter)
{
    int N_full = (int)floor(sqrt(t / (2.0 * M_PI)));
    mpfr_prec_t prec_fast = 64;   /* phase 1 : 1 limbe SIMD */
    mpfr_prec_t prec_full = 116;  /* phase 2 : 2 limbes, convergence 1e-12 */

    /* ... boucle Illinois ... */
    if (i == iter_switch) {
        /* transition : elargir la precision de toutes les variables mpfr */
        mpfr_prec_round(ma, prec_full, MPFR_RNDN);
        /* ... idem mb, mfa, mfb, mc, mfc ... */
    }
    double fc_d = Z_rs_mpfr_ntermes(c_d, N_full, prec_use);
    /* ... mise a jour Illinois classique ... */
}
\end{lstlisting}

Côté Python, \texttt{compute\_zeros\_v6.py} charge la nouvelle signature par ctypes
et appelle \texttt{illinois\_refine\_adaptive} avec \texttt{ITER\_SWITCH = 8},
\texttt{MAX\_ITER = 50}, et un fallback vers \texttt{illinois\_refine} classique en cas
d'exception.

%----------------------------------------------------------------------
\section{Benchmarks mesurés}
%----------------------------------------------------------------------

\subsection{Profil phases v7 --- $T = 100\,000$}

\begin{center}
\begin{tabular}{lrrrr}
\toprule
Phase & Temps cumulé & Appels & ms/appel & \% mur$\times$W \\
\midrule
\texttt{illinois\_C} & 5\,799,66 s & 137\,934 & \textbf{42,05} & 78,2\,\% \\
\texttt{detection}   &    37,80 s  & 4         & 9\,450 & 0,5\,\% \\
\texttt{turing}      &    22,05 s  & 1         & 22\,046 & 0,3\,\% \\
\texttt{mpmath\_petit\_t} & 0,78 s & 138       & 5,65  & 0,0\,\% \\
\bottomrule
\end{tabular}
\end{center}

\subsection{Comparaison v6 vs v7}

\begin{center}
\begin{tabular}{lrrrr}
\toprule
Run & v6 (170 bits) & v7 (64$\to$116 bits) & Gain \\
\midrule
T=5\,000 --- vitesse (z/s)    & 136       & \textbf{1\,808}    & $\times 13$ \\
T=5\,000 --- ms/appel         & 23,50 ms  & \textbf{1,47 ms}   & $\times 16$ \\
T=10\,000 --- vitesse (z/s)   & $\approx$66 & \textbf{408}     & $\times 6{,}2$ \\
T=10\,000 --- ms/appel        & $\approx$50 ms & \textbf{3,71 ms} & $\times 13$ \\
T=100\,000 --- vitesse (z/s)  & $\approx$20 & \textbf{74,49}  & $\times 3{,}7$ \\
T=100\,000 --- ms/appel       & $\approx$130 ms & \textbf{42,05 ms} & $\times 3{,}1$ \\
T=100\,000 --- durée          & $\approx$113 min & \textbf{30,9 min} & $\times 3{,}7$ \\
\midrule
T=100\,000 --- zéros          & 138\,069  & 138\,069           & identique \\
T=100\,000 --- manquants      & 0         & 0                  & identique \\
T=100\,000 --- Turing         & COMPLET   & COMPLET            & identique \\
\bottomrule
\end{tabular}
\end{center}

%----------------------------------------------------------------------
\section{Analyse du gain en fonction de $t$}
%----------------------------------------------------------------------

Le gain décroît avec $T$ car le coût d'un appel à \texttt{Z\_rs\_mpfr\_ntermes}
est $\mathcal{O}(N_\text{full})$, avec $N_\text{full}(t) = \lfloor\sqrt{t/2\pi}\rfloor$ :

\begin{equation}
  N_\text{full}(T = 5\,000)   \approx 20, \quad
  N_\text{full}(T = 10\,000)  \approx 40, \quad
  N_\text{full}(T = 100\,000) \approx 126.
\end{equation}

À $T = 5\,000$ ($N \approx 20$), la boucle \texttt{for n=1..N} est courte et la
différence 64 bits vs 170 bits (routines SIMD vs génériques) domine : gain $\times 16$.
À $T = 100\,000$ ($N \approx 126$), la boucle sur les 126 termes domine et le gain
de précision est dilué : gain $\times 3,1$ sur \texttt{ms/appel}.

La relation empirique est :
\begin{equation}
  \text{gain}(T) \approx \frac{\alpha}{N_\text{full}(T)},
  \quad \alpha \approx 300 \quad (\text{ajusté sur les mesures}).
\end{equation}

\subsection*{Précision finale}

Les deux versions convergent vers le même pseudo-zéro RS, dont la distance au vrai
zéro de Riemann est bornée par le biais structurel de la formule de Riemann--Siegel :
\begin{equation}
  \varepsilon_\text{RS}(t) = \mathcal{O}\!\left(t^{-3/2}\right)
  \approx \begin{cases}
    10^{-4} & t = 300 \\
    10^{-6} & t = 77\,000
  \end{cases}
\end{equation}
Cette précision est largement suffisante pour la validation Turing--Backlund
(espacement minimal entre zéros $\approx 2\pi/\ln(t/2\pi) \approx 0.5$ à $t = 100\,000$).

%----------------------------------------------------------------------
\section{Résultat final v7 --- $T = 100\,000$}
%----------------------------------------------------------------------

\begin{center}
\begin{tabular}{ll}
\toprule
Paramètre & Valeur \\
\midrule
Script         & \texttt{compute\_zeros\_v6.py} (illinois\_refine\_adaptive) \\
Zéros trouvés  & \textbf{138\,069} \\
Manquants      & \textbf{0} \\
Turing--Backlund & \textbf{COMPLET} \\
LMFDB          & 19/20 à $< 10^{-10}$ \\
Durée          & \textbf{30,9 min} \\
Vitesse        & \textbf{74,49 z/s} \\
Gain vs v6     & \textbf{$\times 3.7$} \\
Log            & \texttt{logs/run\_T100k\_v7\_20260611\_1424.log} \\
\bottomrule
\end{tabular}
\end{center}

%----------------------------------------------------------------------
\section{Conclusion et perspectives}
%----------------------------------------------------------------------

La v7 démontre qu'une réduction de précision \emph{bien ciblée} (64 bits pour 8
itérations Illinois sur 15--20) apporte un gain substantiel ($\times 3,7$) sans
aucune perte de correction mathématique.

\subsection*{Limitation connue}

$N_\text{fast} = N_\text{full}/4$ a été écarté car les termes manquants changent
le signe de $Z_\text{RS}$ et brisent la convergence.
Toute tentative future de réduction des termes devra intégrer un \emph{guard de signe}
(comparaison avec \texttt{Z\_double} avant d'accepter le signe de $Z_\text{RS}$ tronquée).

\subsection*{Pistes post-v7}

\begin{enumerate}
  \item \textbf{Termes RS plus rapides} : implémenter $C_2, C_3$ (corrections d'ordre
    supérieur) pour réduire le biais $\varepsilon_\text{RS}$ à $\mathcal{O}(t^{-5/2})$.
  \item \textbf{Vectorisation interne} : batch de plusieurs appels $Z_\text{RS}(c_i)$
    simultanés dans Illinois (SSE/AVX pour la boucle sur $n$).
  \item \textbf{Phase 1 en double pur} : remplacer \texttt{mpfr} à 64 bits par
    \texttt{double} IEEE 754 ($\sim 52$ bits) pour les itérations de réduction rapide,
    évitant l'overhead d'allocation \texttt{mpfr}.
\end{enumerate}

\vspace{1.5em}
\hrule
\vspace{0.5em}
\small\textit{Document produit le 11 juin 2026. Commit \texttt{8637098} sur la branche
\texttt{Riemann\_Lab\_C}. Voir \texttt{JOURNAL.md} du wiki pour l'historique complet.}

\end{document}
